\documentclass[11pt,a4paper]{article}
\usepackage[T1]{fontenc}
\usepackage[utf8]{inputenc}
\usepackage{amsmath,amssymb,xcolor}
\usepackage[margin=25mm]{geometry}
\usepackage[colorlinks=true,urlcolor=blue]{hyperref}
\definecolor{nexorablue}{HTML}{2459B5}
\setlength{\parindent}{0pt}
\setlength{\parskip}{7pt}
\setlength{\emergencystretch}{2em}
\renewcommand{\familydefault}{\sfdefault}
\begin{document}

{\small\bfseries\color{nexorablue} NEXORA CPGE}\par
{\small MP / NEX-STR-01 / Énoncé}\par
{\LARGE\bfseries Un sous-groupe défini par une contrainte additive}\par
\vspace{0.25em}\hrule\vspace{0.6em}

Dans $(\mathbb Z^2,+)$, soit $H=\{(a,b)\mid a-2b\in3\mathbb Z\}$. Montrer que $H$ est un sous-groupe et que tout élément de $H$ s’écrit de manière unique comme combinaison entière de $(3,0)$ et $(2,1)$.


\end{document}
